\documentclass[11pt,a4paper]{article}
\usepackage[T1]{fontenc}
\usepackage{lmodern}
\usepackage[margin=25mm]{geometry}
\usepackage{amsmath,amssymb,amsthm,mathtools}
\usepackage{microtype}
\usepackage{booktabs,tabularx,array,longtable}
\usepackage{xcolor}
\usepackage[colorlinks=true,linkcolor=blue!40!black,citecolor=blue!40!black,
  urlcolor=blue!40!black]{hyperref}
\hypersetup{pdftitle={A Single-Exponential Exponent for Arithmetic Progressions},
  pdfsubject={Quantitative proof modification with an explicit dependency audit},
  pdfauthor={}}
\setlength{\parindent}{1.2em}
\setlength{\parskip}{2pt}
\setlength{\emergencystretch}{2em}
\allowdisplaybreaks[1]
\newtheorem{theorem}{Theorem}[section]
\newtheorem{lemma}[theorem]{Lemma}
\newtheorem{proposition}[theorem]{Proposition}
\theoremstyle{definition}
\newtheorem{definition}[theorem]{Definition}
\newcommand{\M}{\mathcal M}
\newcommand{\E}{\mathbb E}
\newcommand{\Z}{\mathbb Z}
\newcommand{\R}{\mathbb R}
\newcommand{\g}{\mathfrak g}
\newcommand{\h}{\mathfrak h}
\newcommand{\n}{\mathfrak n}
\newcommand{\Ad}{\operatorname{Ad}}
\newcommand{\ev}{\operatorname{ev}}
\newcommand{\Lip}{\operatorname{Lip}}
\newcommand{\wt}{\operatorname{wt}}
\newcommand{\corr}{\operatorname{corr}}
\newcommand{\poly}{\operatorname{poly}}
\newcommand{\im}{\operatorname{image}}
\newcommand{\one}{\mathbf 1}

\begin{document}
\begin{center}
{\Large\bfseries A Single-Exponential Exponent\\[3pt]
for Arithmetic Progressions}\\[7pt]
{\normalsize A quantitative proof modification for OpenAI/math 159}\\[5pt]
{\small October 8, 2026}
\end{center}

\begin{abstract}
This note quantifies the exponent in the density bound of
\emph{Quasipolynomial Bounds for Arithmetic Progressions}
(OpenAI/math, September 23, 2026). Relative to that manuscript's
analytic framework, we obtain
\(\varepsilon_k=\exp(-k^C)\), for an absolute constant \(C\).
The proof keeps a polynomial geometric degree and replaces two
sequential inverse-theorem chains by degree descents. A direct
Cauchy--Schwarz calculation integrates supplied derivative correlators,
while effective nilpotent cube completion removes one vertex in
relative detection. Separating structural dimension from logarithmic
precision controls the remaining branching arguments. The geometric
and analytic modifications and the final numerical schedule are given
explicitly, relative to the source results listed below.
\end{abstract}

\section{The statement and the scope of the proof}\label{sec:result}

For integers \(k\ge3\), let \(r_k(N)\) be the largest cardinality of a
subset of \(\{1,\ldots,N\}\) containing no nonconstant \(k\)-term
arithmetic progression. All logarithms are natural. Write \(\M\) for the
198-page manuscript \cite{openai}; its section and page numbers below
refer to that version.

\begin{theorem}[Quantitative consequence of the source framework]\label{thm:main}
Subject to the analytic results of \(\M\) specified below, there is an
absolute constant \(C>0\) such that for every \(k\ge3\) there are
constants \(K_k,c_k>0\) satisfying
\begin{equation}\label{eq:main}
 r_k(N)\le K_k N\exp\!\left[-c_k(\log N)^{\varepsilon_k}\right],
 \qquad \varepsilon_k=\exp(-k^C),\qquad N\ge3.
\end{equation}
\end{theorem}

\paragraph{Imported results.}
We use the bottom absolute-patch theorem (\(\M\), Lemma 2.2), the
sampler and rank-transfer contracts (Sections 3--4 and Appendices F--H),
the comparison and exact-marking framework (Sections 5--8 and Appendices
A--E and I), the extraction of Section 9, and the density-increment
construction and parameter order of Section 10. We retain their
analytic conclusions, including the prescribed dependencies of warm
bounds. The algebraic modifications and the dependence of the numerical
powers on the fixed orders are the subjects of this note. This is a
proof relative to that manuscript, not a replacement proof or an
independent verification of all its analytic arguments.

The external inverse input is the proof of Leng--Sah--Sawhney
\cite{lss}, using Leng's equidistribution result \cite{leng}. Its
published fixed-order notation does not by itself specify the
order-dependence needed here; Section~\ref{sec:inverse} records the
additional budget calculation. The supplied-derivative integration statement needed for the
improvement is proved in Proposition~\ref{single:integration}.

\paragraph{Budget conventions.}
A growing budget contains dimensions, logarithmic rational heights,
logarithmic Lipschitz bounds, logarithmic inverse errors and correlations,
and logarithmic scale requirements. Order-dependent positive factors
\(K_q\) may be arbitrarily large: they are coefficients, not powers of
the growing budget. We retain every power of that budget and every
successive substitution of one such bound into another. Every exponent
implicit in \(O(1)\) below is absolute. For example, a map
\(x\mapsto K_q(2+x)^b\), applied successively \(m\) times, has
sufficient degree \(b^m\); calling \(m\) fixed for a fixed \(q\)
does not justify dropping it from an estimate in \(q\).

Set
\[
 s=k-2,\qquad D=256k^2,\qquad p=2+\log(1/\alpha)
\]
for a starting density \(\alpha\). The original multiplicative density
gain requires \(T\le K_kp\) rounds. At weight \(h\), the determining
data consist of integer coordinates \(b_h\), a rational value space
\(W_h\), and targets
\[
 C_h(u,b_{<h})\in W_h,\qquad \deg_{\wt} C_h\le h.
\]
The spatial variable has weight one and \(b_h\) has weight \(h\).
The paired determining cutoffs force a unique padded integer tuple.
Writing \(Q_h=\log(2/w_h)\), a bound is \emph{warm at \(h\)} if it
may use dimensions, \(p\), and already chosen data at weights strictly
above \(h\), but none of the incoming \(Q_i\) for \(i\le h\).
A preliminary bound at \(h\) may also use \(Q_h\). Once retained as
history, such data are warm at every subsequent layer \(j<h\).

The principal bounds proved below are
\[
 D=O(k^2),\qquad
 \max_{r\le D+2}I_r,\ \max_{q\le D}R_q,\ F,B
       \le\exp(D^{O(1)}).
\]
Here \(I_r\) is an inverse-theorem budget exponent, \(R_q\) is the
relative-detector exponent of Section~\ref{sec:relative}, and \(F\) controls the
fresh absolute rank, and \(B\) controls the source's local polynomials
and their descending numerical schedule. The last density calculation
uses the original rank ledger.

\section{A persistent record with quadratic degree}\label{sec:macro}

A completed active return has two kinds of controlled factors: a slow
factor in normalized coordinates and a rational factor in integer
coordinates. Keeping these factors on their original sides allows the
middle orbit to retain its original degree.

For a filtered Lie group, call a polynomial \emph{jointly filtered} when
each monomial of weighted degree \(a\) in its logarithm has coefficient
in the \(a\)-th Lie layer. Constants lie in the first layer. Integer
and normalized labels have weight one. A \emph{marking} is the fixed
projection \(\pi:G\to Y\) recording the determining equations.

\begin{lemma}[Ordered active return]\label{lem:active}
In the setting of Theorem 8.2 of \(\M\), suppose the input orbit is
jointly filtered in all actual parameters and its marking is a fixed
filtered polynomial \(P_Y(X)\) of the site variables. The group and
marking may themselves be finite charged families in older dummy
variables, which are coefficients rather than actual arguments.
The return can be chosen in the form below, with each of \(E,C,R\)
jointly filtered in the original filtration:
\begin{equation}\label{eq:ECR}
 g^\sharp(X,\theta,\widetilde t,z)
   =E(X,\theta,z)\,C(X,z)\,R(X,\widetilde t,z),
 \qquad
 \pi E=\pi R=1,\quad \pi C=P_Y(X).
\end{equation}
Here \(X\) collects the site variables, \(z\) the short labels, and
\(\theta,\widetilde t\) the normalized and integer long labels.
The factor \(C\) has no long-label argument; \(E\) has controlled coefficients
in normalized site, residual, and label coordinates; and \(R\) has
controlled common denominator in raw integer coordinates. No
coefficient bound is required for \(C\). The original score slice,
both expected-payoff representative identities, and the field
estimates are retained.
\end{lemma}

The additional induction proving the lemma is given in
Section~\ref{sec:active}. We first show how its format closes the macro
argument.

\begin{lemma}[Transport of entire residual features]\label{lem:transport}
At successive active layers, a historical normalized real label can be
stored as an affine expression in the current normalized spatial
coordinates and named higher residuals. At final reassembly every such
feature has weighted degree at most \(s\). Unfrozen historical short
integer labels have a finite menu of whole site functions, with
preliminary cost at their creation layer and warm cost at every lower
layer. These statements concern the exactly marked payoff and retain
both of its zero cases. Only active layers create these histories.
Old determining slots of weight \(h>s\) remain in the canonical
passive site multipliers and are not included among these features.
\end{lemma}

\begin{proof}
Fix one edge at layer \(j\), and write its parent higher raw target as
\(C_h^P(u,b_j,\ldots,b_{h-1})\), where \(j<h\le s\); lower variables already
fixed at this parent have been absorbed into \(u\). The edge has an
affine spatial map \(\psi(t)\) and an exact integer current lift
\(B_j^\psi(t)=\beta_j^\psi(t)+b_{j,0}\). By the definition of the
stopped tree in \(\M\), Section 3.5, p.~23, the child targets are
defined by the formal identity
\begin{equation}\label{eq:target-transport}
 C_h^\psi(t,b_{j+1},\ldots,b_{h-1})
 =C_h^P(\psi(t),B_j^\psi(t),b_{j+1},\ldots,b_{h-1}).
\end{equation}
Undo any fixed triangular determining reindexings on both sides, using
their recorded exact inverses. For the named raw residuals
\(\rho_h=b_h-C_h-\ell_h\), it follows that
\begin{equation}\label{eq:residual-transport}
 \rho_h^\psi(t,b_{>j})
 =\rho_h^P(\psi(t),B_j^\psi(t),b_{>j}).
\end{equation}
The centered residual used by the sampler is an affine translate and
sign change of this feature:
\(y_h=-\rho_h+(b_{h,0}-c_h-\ell_h)\).

Use the fixed prepared offsets \(\ell_h\) in these residuals;
newly selected smaller-box centers remain cutoff parameters. The
canonical tuple is defined on one fixed prepared padded chart:
take its true current raw lift, then the successive higher padded
lifts, using the prepared offsets rather than path-dependent new
centers. Outside that chart use the common zero extension.
In the next external field, replace an old residual \emph{as an entire
feature} by its parent version evaluated at this canonical tuple.
On actual paths this agrees by \eqref{eq:residual-transport}. On an
exactly marked return, the determining transformation is the parent
transformation, and a contributing determining cutoff forces that same
canonical tuple. Keep the whole integer tuple sum. Each inadmissible
tuple contributes zero to both expressions; for an admissible tuple
all substituted features agree. If no padded tuple exists both payoffs
are zero. This proves separate equality for the two unsigned payoffs.
Off the scoring locus this is a definition of the external field, just
as in \(\M\), Lemma 7.5; no claim that rounding is a polynomial is made.

During the next micro return, the previous spatial coordinates are
parameter labels. After the current macro extraction and freezing of
nonpivot parameters, \(\M\), (9.4), expresses these normalized real
labels affinely in \(u/H\), \(y_j/v_j\), and the normalized active
sheet coordinates. The latter are linear in the same residual frame.
The normalized extra inactive labels are also affine in that frame
after their product columns have been fixed, by (9.2). Update only
these affine spatial features, and transport the already named higher
residuals by \eqref{eq:residual-transport}. This proves the affine
feature invariant by induction. At final reassembly insert
\(\rho_h=b_h-C_h(u,b_{<h})-\ell_h\); it has weighted degree at most
\(h\le s\). In particular, one never composes an arbitrary degree-\(h\)
child target directly with a degree-\(j\) pivot, which could have
degree \(hj\).

For the integer menu, an unfrozen extra short label created at \(h\)
has the form \(w(u)=\lambda\beta_h^{\rm can}(u)+c\), with integral
\(\lambda,c\). Lemma 3.1 and (3.1) of \(\M\) give
\[
 \|\lambda\|\le \exp(\poly(d_h))K_i,
 \qquad K_i\le k_i/v_h.
\]
For a short axis, \(k_i\), \(v_h^{-1}\), and the product-column
offset have preliminary bounds. Keep the row in the original raw
coordinates by undoing representation reindexings. If one writes
\(\beta_h=b_h-b_{h,0}\), the additional offset is controlled as well:
\(\lambda\) annihilates \(W_h\), \(c_h\in W_h\), and the
feasible-center relation (3.6) controls \(b_{h,0}-c_h-\ell_h\).
For \(m_h\) labels and \(\|\lambda\|,|c|\le R_h\), there are at most
\[
 (2R_h+1)^{m_h(d_h+1)}
\]
row-and-offset lists. Their logarithmic count is polynomial in the
structural dimensions and the preliminary budget. These lists specify
\emph{whole functions} at the fixed later parent. No regularity or
coefficient bound on the canonical site map is used. Annihilation of
\(W_h\) also bounds the values throughout its padded chart. All these
costs are warm at \(j<h\).

For physical reassembly replace canonical integers only in these
external features by the formal summation integers. The equality just
proved holds on the full chart and exact-inactive scoring locus. This
is the locus on which the two representative identities are required
in (9.8) of \(\M\); apply them before freezing stochastic constants.
\end{proof}

\begin{proposition}[Persistent macro record]\label{prop:macro}
Starting from an absolute patch of degree \(s\) and rank \(d_0\), the
macro argument can retain the original \(d_0\) absolute slots and a
middle orbit of joint degree at most \(s\). Its final physical formula
has absolute-weight ceiling at most \(s(s+1)\). The representation,
field, score, and extraction contracts of \(\M\) are preserved, with
warm costs at each later layer. The padded ceiling \(D=256k^2\)
therefore suffices.
\end{proposition}

\begin{proof}
\textit{The retained data.}
Store the return as \(E_{\rm hist}CR_{\rm hist}\), and use only \(C\)
as the next active orbit. Historical integer short labels remain
formal coefficient variables, with \emph{one total-charge cap} \(s\)
over all their blocks. Raw and normalized copies of a short label
always use the same variable and its fixed normalizing divisor.
Historical long labels remain in the outer factors. A deep
coefficient lattice evaluates into the pointwise lattice at every
integer short label; arbitrary real labels have no such property.
Before the next positive-filtration recursion, normalize any
degree-zero family constant as in the proof of
Lemma~\ref{lem:active}.

A newly promoted determining coordinate retains its original weight
in \(C\). In outer-factor coefficients use its canonical site value:
the left factor uses the bounded normalized residual frame, and the
right factor uses raw integer coordinates. Exact marking identifies
these values at every tested point. Fix the integral determining
reindexings before recording history, and keep their exact
triangular inverses. They preserve the original determining degrees.
Later congruence restrictions have their creation-layer costs, hence
warm costs at every lower layer.

\textit{Ordered representation.}
Normalize the middle factor only. For an integral triangular
reindexing \(T\) of the absolute tuple, the correct expression is
\begin{equation}\label{eq:reindex}
 E_{\rm hist}(CT)(T^{-1}R_{\rm hist}T).
\end{equation}
This is the old physical map followed by \(T\), which permutes the
complete integer tuple sum. Conjugation by the integral map preserves
controlled rational denominators; its numerator need not be bounded.
Coefficientwise reduction gives \(CT=Sg\), where \(S\) is the bounded
full-weight map. It is independent of the historical variables:
their positive degree excludes them from a full-weight slot term.
If \(g=E_{\rm new}C_{\rm new}R_{\rm new}\), collect the result as
\begin{equation}\label{eq:collect}
 \bigl(E_{\rm hist}\Ad_S E_{\rm new}\bigr)
 (SC_{\rm new})
 \bigl(R_{\rm new}R'_{\rm hist}\bigr),
 \qquad R'_{\rm hist}=T^{-1}R_{\rm hist}T.
\end{equation}
Only bounded maps conjugate slow factors. Rational right factors
stay on the right. The triangular subgroup in which determining
transformations have no absolute arguments is bracket closed;
forgetting the absolute slots is strong, since a determining shear
lifts while fixing them. Thus this normalization preserves the
exact determining projection.

\textit{The quotient fields.}
At site \(u\), define the field on the coefficient-family quotient by
the original cutoff sum evaluated on
\begin{equation}\label{eq:field}
 E_{\rm site}(u)\,S\,
       \ev_{w(u)}(Y)\,R_{\rm site}(u).
\end{equation}
Canonical site substitution occurs only in the external coefficients
and in \(\ev_{w(u)}\). All original determining and absolute integer
variables are still summed. On an exactly marked orbit, any nonzero
determining cutoff forces the unique canonical tuple. Hence both
payoffs agree exactly with the originals, including their zero
cases.

A common deep coefficient lattice, further deepened for the rational
right factors, makes \eqref{eq:field} well defined. On bounded group
representatives, the left map and \(S\) are bounded; the right map
can be reduced to a bounded right-coset representative determined by
its denominator and coefficient residues. The triangular inverse
in the representation argument then bounds all contributing integer
tuples and derivatives. Bounded-range integer evaluation has
controlled Lipschitz norm. This gives the required Lipschitz bound,
independently of the unrestricted coefficients of \(C\).
Triangular uniqueness preserves the cap one.

\textit{The two field meshes.}
After extraction, product-column labels are frozen. The only
historical short labels that remain variable are integral affine
maps from extra inactive rows. Those rows annihilate the prepared
value space. The chart estimate has the form
\[
 \|z_i\|\le \exp(\poly(d))K_i,\qquad K_i\le k_i/v,
\]
where \(v\) is the chart width. For a short axis, its scale \(k_i\),
the integral offset, and its values throughout the fixed padded
chart have preliminary bounds. The factor \(v^{-1}\) is included
in that bound. At every later layer all these costs are warm.
The fixed determining polynomials therefore turn the finite
row-and-offset menu into a finite menu of entire functions
\(u\mapsto w(u)\).

Transport historical real labels by Lemma~\ref{lem:transport}, keeping
each higher residual intact until physical reassembly. They are affine
expressions in the current normalized spatial coordinates and named
higher residuals. They are bounded on the same padded chart, even when
they leave a formerly selected unit interval. Retain the original
smooth cutoffs and score slice; no new moving sharp indicator is
inserted. Outside this fixed padded chart use the zero extension
from the representation construction of \(\M\).

At one site, mesh bounded left values, full-weight coefficients,
bounded integer evaluations, right cosets, and inherited cutoff
parameters. For a mesh uniform over all sites, mesh the bounded
left \emph{coefficient vectors}, the finite whole-function short
menu, rational coefficient residues, and the inherited field
parameters. The canonical site functions are fixed throughout.
These give the distinct point-value and full-field meshes required
by Definition 8.1 of \(\M\).

First pigeonhole the finite tuple of actual historical denominators
and cover moduli, then take their product and a fixed deepening.
Its logarithm is a sum of warm budgets. Taking the least common
multiple of all integers below an upper bound is unnecessary.
After product-column freezing, historical raw labels are affine
in raw spatial and determining integers and the remaining extra
inactive labels. Later integer sections may have arbitrarily large coefficients;
they change only the rational numerators, which are recorded
modulo the common deep grid. The same divisibility gives the
long-label periods. Selection pays a warm loss in the fraction of
available productive paths; every retained path keeps its original
full score and slice.

These history costs may depend on the creation-layer width and are
only preliminary there. They are not inserted into that layer's warm
relative width loss. The boundary estimates of \(\M\), pp.~71--72,
use the cap one for arbitrary unsigned payoffs; the averaging on
pp.~73--74 uses the total selection measure; and domain localization
on pp.~74--75 does not approximate the polynomial payoff. Its
Lipschitz bounds apply to the normalized pivot and interval tests.
In particular, the micro proof's preliminary short threshold is
distinct from the warm small-axis cutoff used for exact domain
localization. Larger axes use their normalized affine interval tests.

\textit{Extraction and the final degree.}
For a macro comparison, reassemble the physical formula. On its
scored locus, replace canonical integers in external coefficients
by their corresponding integer summation variables. Exact marking
and the original cutoffs identify the expressions. Current
auxiliary labels now enter only through this polynomial formula,
and both expected representative identities apply before the
stochastic constants are frozen. Section 9 of \(\M\) uses these
identities, the original slice and masks, cap one, exact
determining and inactive equations, and a degree ceiling. It
does not require bounds on coefficients or spatial derivatives
of the returned physical polynomial. Its extraction and
localization therefore apply with the original payoffs.

Every recorded historical-label substitution has final
spatial/determining weighted degree at most \(s\). For a monomial
of original absolute-input weight \(h'\) and other charge \(d\),
the joint invariant gives \(h'+d\le h\), where \(h\le s\) is
its absolute destination weight. Give each absolute slot the
temporary weight \((s+1)h\). Simultaneous substitution gives
\begin{equation}\label{eq:quadratic}
 (s+1)h'+sd\le (s+1)(h'+d)\le(s+1)h.
\end{equation}
This reweights absolute inputs and destinations together.
Consequently the physical ceiling is \(s(s+1)\).

These temporary weights are used for extraction and final
compression. The next middle orbit still uses the original
weights, with determining slots and historical integer variables
formal; its outer factors again enter the fields. This closes
the induction without repeated degree inflation. No family or
partner coordinate becomes an absolute slot, so the root rank
is still \(d_0\), independent of old numerical widths.

Finally, \(m\) historical variables of total charge at most \(s\)
have at most
\[
 \binom{m+s}{s}\le(m+1)^s
\]
monomials per scalar coordinate. There are at most \(s\) macro
layers and \(s\) internal degree reductions per return.
Rational systems, coefficient arrays, and finite meshes are
handled simultaneously. The added operations have polynomial
sequential depth in \(s\); their quantitative bounds are included
in Section~\ref{sec:numerical}.
\end{proof}

\section{The charged active return}\label{sec:active}

We prove Lemma~\ref{lem:active} by supplementing the degree induction
of Theorem 8.2 of \(\M\). The point is to retain total degree when
short variables become coefficient variables, while each recursive
child still has strictly smaller positive degree.

\begin{proof}[Proof of Lemma~\ref{lem:active}]
\textit{Charged families and positive normalization.}
Record rational coefficient spaces \(A_{d,\ell}\), decreasing in both
indices, with
\[
 [A_{d,\ell},A_{e,m}]\subset A_{d+e,\ell+m},\qquad
 A_{0,\ell}=A_{1,\ell},\qquad A_{r+1,\ell}=0,
\]
and one finite total-charge cutoff in \(\ell\).
Initially \(A_{d,\ell}=\g_{\max(d,\ell)}\), with
\(\g_0=\g_1\). Promoting a short block to a dummy \(w\) replaces
the array by
\begin{equation}\label{eq:family}
 A'_{d,\ell}
   =\left\{\sum_\beta w^\beta a_\beta:
                 a_\beta\in A_{d,\ell+|\beta|}\right\}.
\end{equation}
The positive layers are \(\exp A'_{i,i}\), and the
degree-zero group is \(\exp A'_{0,0}\). Quotients use images
of these spaces, and refinements use the actual intersections
before taking images. Every raw and normalized occurrence of a
promoted label is changed simultaneously, including those in
marks and ancestor factors.

Perform the common-site marking correction before fast-algebra
selection. A coefficientwise ordinary filtered section commutes
with evaluation at the actual short label, so the inserted
correction evaluates to the identity. Normalize the fixed marked
constant by a bounded left factor and a rational right factor.
The normalized mark takes values in the positive marked group.
Let \(c=g(0,0;w)\) be the constant family element after this
marked normalization.
The remaining constant coset modulo the positive group can now
be removed in the kernel: choose a positive-layer lift \(c_1\)
of its marked image, and split the difference
\(n_0=cc_1^{-1}=e_N\gamma_N\) into a bounded kernel element
and a kernel lattice element. Replace the family orbit \(g\) by
\(e_N^{-1}g\gamma_N^{-1}\), retaining both factors for restoration.
This lies in the positive group because its constant coset
does. These are family-group operations, so all charge bounds
are retained. A deep source lattice with the prescribed image
handles the marked rational factor; its numerator need not be
bounded.

\textit{Refiltration with the old degree retained.}
The degree-zero action is made compatible with the refiltration
by adjoining to the detected fast algebra the rational graded
ideal
\[
 S_i=\im\!\left(A'_{i+1,i}\longrightarrow
                    A'_{i,i}/A'_{i+1,i+1}\right).
\]
The bracket rule makes degree-zero conjugation trivial modulo
\(S_i\), and \(S_r=0\). Thus saturation preserves the detected
top and its injective marked projection, while the resulting
refiltration is normalized by the degree-zero group.

Write \(G_i\) for the old positive layers and \(H_i\) for
the refiltration from Lemma 7.6 of \(\M\). The prescribed
full factors preserve old joint degree: a grade-\(i\) symbol
correction is lifted into \(G_i\); the off-space substitution
preserves weighted degree; and the constant adjustment sees
only old grade one, the higher short degrees already lying
in \(G_2\subset H_1\). A pure-long grade-\(i\) symbol is independent
of the short variables, since every coefficient of positive
short degree lies in old \(G_{i+1}\); its kernel lift therefore
preserves total degree. Filtered multiplication and conjugation
preserve these bounds. The refiltered input therefore satisfies
\begin{equation}\label{eq:mixed}
 [v^\alpha z^\beta]\log h
       \in \h_{|\alpha|}\cap\g_{|\alpha|+|\beta|},
 \qquad \h_0=\h_1.
\end{equation}
Here \(v\) denotes the long variables and \(z\) the current
short variables.

Since \(G_{i+1}\subset H_i\subset G_i\), the layers
\[
 H_1\supset G_2\supset H_2\supset G_3\supset\cdots
\]
form one rational flag, with its strong image in the marked
group. A rational section adapted to this flag preserves the
required intersections. It gives a fixed exact lift \(p(X,z)\)
of the refiltered mark with
\[
 [X^\gamma z^\beta]\log p\in
       \h_d\cap\g_{d+|\beta|},\qquad d=\wt\gamma.
\]
Only this interleaved local flag is used; historical charge is
already encoded in the coefficient-family group.

\textit{One lower-degree child.}
Let \(q:H\to Q=H/H_r\), and let \(N\) be the actual
kernel of the marking on \(H\). Top injectivity identifies
\(N\) with the kernel of the quotient marking via \(q|_N\).
Promote the current short block to a new dummy \(w\) and form
\begin{align*}
 \mathcal Q_i
   &=\left\{\sum_\beta w^\beta a_\beta:
           a_\beta\in q(\h_i\cap\g_{i+|\beta|})\right\},\\
 \mathcal N_i
   &=\left\{\sum_\beta w^\beta n_\beta:
           n_\beta\in q(\n\cap\h_i\cap\g_{i+|\beta|})\right\}.
\end{align*}
At index zero use \(\h_0=\h_1\) and \(\g_0=\g_1\).
These are images of the displayed intersections.
The bracket relation
\([\mathcal Q_i,\mathcal N_j]\subset\mathcal N_{i+j}\)
defines the split rational family
\(\mathcal N\rtimes\mathcal Q\), marked by its second
coordinate. Its canonical section is \(a\mapsto(1,a)\).
The map \((n,a)\mapsto(na,a)\) embeds it into two copies of the
quotient family, so it is a rational nilpotent group.
Both degree-\(r\) spaces vanish, so its positive group has
degree at most \(r-1\).

The local pair
\[
 \bigl(qh\,(qp)^{-1},\,qp\bigr)
\]
has the required kernel coefficients by \eqref{eq:mixed}.
Correct its marking to the common site with the canonical
section, and apply the degree-zero normalization above.
The resulting input is jointly filtered in all actual
variables: substituting a weight-\(d\) site monomial can
produce actual degree \(a\le d\), and a dummy degree \(b\)
is charged in the old layer \(\g_{d+b}\subset\g_{a+b}\).
Apply the induction to this single child, leaving \(w\)
unevaluated.

Restore the child's preliminary degree-zero factors, so its
exact second coordinate is \(qp(X,w)\). Multiplying by the
inverse of its canonical marked section isolates the kernel
family. Lift this family through \((q|_N)^{-1}\), multiply
by \(p(X,w)\), and evaluate \(w=z\). Its degree-\(a\),
dummy-degree-\(b\) coefficients lift into
\(\n\cap\h_a\cap\g_{a+b}\). Thus reconstruction restores
the original joint degree exactly, including synchronized
normalized copies. A common deep coefficient lattice
evaluates inside the pointwise lattice at every integer
\(z\). Fields pull back through the homomorphism
\((n,a)\mapsto na\); pullback preserves supremum meshes,
and bounded short evaluation has controlled Lipschitz cost.
No metric bound on the particular marked lift \(p\) is used.

\textit{Ordered factors and the original score.}
Keep the child return in the full split marked group.
The joint-fiber reconstruction \(T\) is equivariant for
actual kernel factors:
\[
 T\bigl((e,1)c(r,1)\bigr)=\widetilde e\,T(c)\,\widetilde r.
\]
Indeed the two sides have the same quotient and exact
mark, and the joint embedding is injective. The lifts
\(\widetilde e,\widetilde r\) are rational and filtered.
Collect them through the prescribed bounded left and
rational right factors of \(\M\). For the local notation
\(K_E\rho(E_F)\) and \(\rho(R_F)K_R\), this gives
\[
 E=K_E\Ad_{\rho(E_F)}(\widetilde e),\quad
 C=\rho(E_F)T(c)\rho(R_F),\quad
 R=\Ad_{\rho(R_F)^{-1}}(\widetilde r)K_R.
\]
The left conjugation has bounded normalized coefficients,
and the right one has controlled rational denominator.
The middle has no long-label argument.

If \(L\) is the marking restorer of \(\M\), use
\(E(LC)R\). It has the required marking everywhere and
agrees with \(LECR\) on the full exact-inactive forecasting
locus, where \(L=1\). That is the locus of both original
payoffs and their representative identities. At an earlier
central quotient, lift the outer factors separately in
the actual kernel and lift the middle separately.
Any difference from the original selected lift is removed
by the independent central Haar average. Restore the
initial degree-zero factors and project away partners.
The marking and joint degree are exact, and the expected
payoffs, original slice, and additive analytic errors
are those of \(\M\).

At degree zero the scalar-comparison base case applies.
Every positive step has one strictly smaller-degree child.
This proves the lemma. All additional operations are
simultaneous rational linear algebra, bounded-degree
coefficient constructions, and one split product per step.
\end{proof}


\section{The external inverse budget}\label{sec:inverse}\label{sec:numerical}

For degree \(r\), let \(I_r\) be an admissible exponent such that a
\(U^{r+1}\)-norm lower bound \(\delta\) for a one-bounded
\(f:[N]\to\mathbb C\) supplies a correlating
degree-\(r\) nilsequence with
\begin{equation}\label{eq:inverse-budget}
 d+\log M+\log(2+\|F\|_{\Lip})+\log(1/\corr)
 \le K_r\bigl(2+\log(1/\delta)\bigr)^{I_r}.
\end{equation}
During constructions of vector-valued nilcharacters we include the
\emph{logarithm} of their output dimension in the budget. It need not
be the dimension of a nilmanifold.

\subsection{Uniform coordinate and counting operations}

\begin{lemma}[Bulk operations]\label{lem:bulk}
Let \(q\) bound the nilpotence degree, filtration nesting length,
total multidegrees, and orbit-variable count in an operation of
\cite{leng,lss}. Put \(X=2+d+\Lambda\), where \(d\) includes all
growing manifold-coordinate and Lie-generator counts. Vector-output
dimensions and numbers of Fourier alternatives and charts enter
only through their logarithms in \(\Lambda\). The latter also includes
the other logarithmic numerical data just specified. The nonrecursive
rational, coordinate, metric, Fourier, and universal-group operations
used below admit bounds
\[
 X_{\rm out}\le K_q X^{H_q},\qquad H_q\le\exp(q^{O(1)}).
\]
Joint linear algebra and the fixed-grid lattice constructions are
treated as single operations. The bound does not charge the sizes of
unrestricted polynomial-orbit coefficients.
\end{lemma}

\begin{proof}
For an \(m\)-by-\(n\) rational matrix of logarithmic height \(L\),
clear the denominators jointly, at logarithmic cost at most \(mnL\).
Maximal minors and Cramer's rule bound kernels, images, sections, and
intersections by a fixed polynomial in \(m+n\) times \(1+L\).
Stack simultaneous conditions before elimination. A basis of a nested
flag is selected from its spanning lists, followed by one change of
basis; its height need not be recursively powered once per dimension.

The BCH formula has bracket length at most \(q\). A coordinate
expansion uses at most \(K_q(1+d)^{O(q)}\) indexed terms. A triangular
inverse through at most \(q\) nesting blocks has degree at most
\(q^q\). Thus even expanding every needed coefficient array gives
at most
\[
 K_q(1+d)^{\exp(O(q\log(q+1)))}
\]
entries. Denominator clearing, differentiation, coefficient extraction,
and rational elimination on these arrays have the asserted bound.
The same coordinate polynomials compare the metrics on bounded
fundamental representatives. Their Lipschitz constants have logarithm
bounded by \(K_q(1+d)^{\exp(q^{O(1)})}(1+\Lambda)\).
Here one must not power a budget along an arbitrarily long path in
the group metric. There is a direct bound for its entire product.
Write its increments as \(\exp(Y_i)\) and put
\(S=\sum_i\|Y_i\|\). In the tensor algebra truncated at degree
\(q\), the degree-\(a\) coefficient mass of the product of these
exponentials, and then of its logarithm, is bounded by \(K_qS^a\).
Evaluating the Lie words gives
\[
 \left\|\log\prod_i\exp(Y_i)\right\|
 \le K_qM^{K_q}(1+d)^{O(q)}\sum_{a=1}^q S^a.
\]
Coordinate conversion bounds \(S\) polynomially in the total
coordinate path length with an order-controlled degree. Taking the
infimum over paths yields the metric-to-coordinate bound at once.
Polynomial difference factorization gives the converse bound and
weak left-translation estimates. Thus the number of path segments
does not enter the degree of a budget polynomial.

For a rational subgroup, define its Lie algebra by one integer
matrix and use the common weak-lattice denominator. A bounded
lattice point close enough to this rational subspace must lie in
it: a nonzero defining linear form has magnitude at least the
reciprocal common denominator. Its separation threshold has a
logarithmic bound of the same type. Below that threshold the
quotient displacement belongs to the subgroup lattice; above it,
use bounded diameter divided by the threshold. This proves the
subgroup and quotient metric comparisons without dimension descent.

Two apparently longer lattice loops deserve explicit mention.
In \cite[Lemma B.11]{leng}, one first fixes a weak-lattice integer
\(Q_0\), and all the resulting basis coefficients have height at
most \(Q_0^2\). In \cite[Lemma 10.4]{lss}, the common divisibility
constants and weak-lattice inclusions are fixed before collecting
the words. Arbitrarily many collections therefore do not create
arbitrarily many budget powers. Rational powers and periodicity use
the same fixed nesting blocks and joint coefficient denominators
\cite[Lemmas B.13--B.14]{leng}.

A free nilpotent group on \(m\) generators of positive weights and
allowed total weight at most \(q\) has a spanning word list of size
at most
\[
 \sum_{a=1}^q 4^a m^a\le K_q(1+m)^q.
\]
Relations within a fixed word type change order-dependent
coefficients; the growing generator count is explicitly retained
in this dimension estimate. The rational operations just proved
then apply to the universal constructions in Section 10 and
Lemma C.6 of \cite{lss}.

On a \(d\)-dimensional torus, Fourier truncation at Lipschitz
constant \(L\) and error \(\eta\) has logarithmic term count
\(O(d^2(\log(2+d)+\log(2+L)+\log(1/\eta)))\).
Finite chart partitions have the corresponding logarithmic count.
The common-function step in \cite[Lemma 6.2]{lss} selects a member
of such a fixed partition; it does not enumerate a net of all
Lipschitz functions. Structure-constant selection costs a polynomial
in \(d\) and \(\log M\). Tensor products add logarithmic output
dimensions. Neither these coordinates nor all Fourier alternatives
are made into separate nilmanifold factors.

The partition construction in \cite[Lemma B.3]{lss} has one child
on the quotient by the last degree layer. At a fixed coordinate
precision \(\eta\), the fibre chart counts have total logarithm
\(O(d\log(C/\eta))\). Local metric comparisons and the product
rule give their Lipschitz bounds; there are at most \(q\) degree
levels. A bump of width \(\eta\) is assigned the valid derivative
bound \(O(\eta^{-1})\). Nilcharacter embedding (Lemma B.4) uses
this partition once. The shear construction (the deferred proof
of Lemma 11.3) has rational word and clock-module bases. Clock
derivatives copy or kill a word rather than lengthening it; its
function is an old function times cutoff factors. The product
rule charges their number and the logarithm of the seam width.
No iteration in that width is required.

This also accounts for the rational quotient and pullback helpers
3.10 and 10.5--10.12 of \cite{lss}, and its degree-rank
factorization in Appendix B. Their homomorphisms are linear in
first-kind coordinates; kernels, images, ideals, and the required
changes of basis are the simultaneous matrix operations above.
The coefficient and scalar helpers in \cite[Appendix A]{leng}
use the same operations. Its scalar Weyl, Vinogradov, and
extrapolation estimates have losses of the form
\(\eta^{-C_q}\) or \(Q^{C_q}\); after logarithms their unspecified
scalar constants are coefficients. They are not powers of the
growing logarithmic budget. Its refined bracket estimates have
explicit factors polynomial in the growing dimensions after
logarithms. The Fourier step in C.6 of \cite{lss} can use the
torus truncation of \cite[Lemma A.5]{leng} in the controlled
local coordinates just described.

The additive-combinatorial primitive \cite[Lemma A.1]{lss} already
has absolute polynomial-logarithmic powers, including its growing
dimension. It can therefore be inserted as one bulk operation. Its
internal character-selection count need not be interpreted as a
sequence of new generic budget powers. Combining this finite list
of coordinate and counting constructions proves the lemma.
\end{proof}

\subsection{Affine numerical dependence in equidistribution}

\begin{lemma}[The complete equidistribution operation]\label{single:equid}
Let \(q\) bound the degree, central step, filtration nesting and
number of orbit variables in Leng's Theorem~4. Let \(d\) be the
input Lie dimension and let \(\Lambda\ge2\) contain the logarithmic
structural, Lipschitz, rationality and inverse-correlation data, together
with relative scale losses and requested side cutoffs. It does not
include the unrestricted orbit coefficients or \(\log N\).
The logarithmic output data and sufficient side requirements are at most
\[
 K_q(1+d)^{E_q}\Lambda,\qquad E_q\le\exp(q^{O(1)}).
\]
The groups supplied to its recursive calls have dimension at most
\(K_qd\), independently of \(\Lambda\). Thus the complete theorem
is one combined-budget operation of degree \(\exp(q^{O(1)})\).
\end{lemma}

\begin{proof}
We refine Lemma~\ref{lem:bulk} by separating \(d\) and
\(\Lambda\). Joint denominator clearing, minors and rational
sections have logarithmic height linear in the original logarithmic
height. Coordinate arrays have at most
\(K_q(1+d)^{\exp(q^{O(1)})}\) entries, and their coefficients
are sums of products of order-bounded length. The entire-product
metric estimate and the fixed-lattice separation argument used above
retain this linear dependence. Fourier selection costs the logarithm
of its number of alternatives, which is a structural polynomial times
the logarithmic Lipschitz and error data.

The scalar Weyl, Vinogradov and extrapolation estimates in
\cite[Appendix A]{leng} use order-only powers of the numerical
heights, errors and retained-scale reciprocals. The refined bracket
bounds in Section~3 have factors of the form
\((\delta/(d^{O(d)}M))^{-O(1)}\); after logarithms, their
\(d\log d\) term is a structural polynomial. Their vector
versions have \(O_q(d)\) coordinates. Simultaneous character
factorization, rationality pigeonholing and the period and progression
operations in Appendices A--B have the same affine dependence.
Consequently a nonrecursive stage has the bound
\[
 \Lambda_{\rm out}\le K_q(1+d)^{H_q}\Lambda,
 \qquad H_q\le\exp(q^{O(1)}).
\]
Its coefficient arrays are charged in \(H_q\). They are not the
Lie dimension passed to the next theorem call: the actual groups are
squares, subgroups and quotients. A square has dimension at most
\(2d\), and at most \(O(q)\) squares occur on a recursive path.

The two calls in Case~2 of \cite[Section~4]{leng} are sequential,
so they must both be retained. Write \(e(a,b)\) for the structural
power at polynomial degree \(a\) and central step \(b\), and
\(e_1(a,b)\) for Case~1. Case~1 has one smaller-central-step
child and the simultaneous recovery operations just listed, hence
\[
 e_1(a,b)\le e(a,b-1)+C H_q.
\]
Case~2 first invokes the theorem on a normalized square of degree
\(a-1\). It then refilters a subgroup of the original group,
selects a rational progression, and invokes Case~1 with the updated
numerical data. The latter input still has Lie dimension at most
the original \(d\). Its two numerical bounds compose as
\[
 \Lambda_1\le K_q(1+d)^{e(a-1,b)+CH_q}\Lambda,
 \quad
 \Lambda_2\le K_q(1+d)^{e_1(a,b)}\Lambda_1.
\]
Thus the structural powers add:
\[
 e(a,b)\le e(a-1,b)+e(a,b-1)+C'H_q.
\]
The abelian base has power \(O(H_q)\). A Pascal majorant gives
\(e(a,b)\le C''H_q2^{a+b+2}\) for \(a,b\le q\), proving
the result. In particular, no power of \(\Lambda\) has been
composed once per node of this branching calculation.
\end{proof}

\subsection{The non-inverse dependency graph}

\begin{lemma}[One order of the original inverse proof]\label{lem:external-body}
With its two lower-order inverse calls removed, an order-\(r\)
construction in \cite{lss} has logarithmic-budget degree
\[
 F_r\le\exp(r^{O(1)}).
\]
\end{lemma}

\begin{proof}
The operation count is by recursive indices, not by the number of
statements with fixed-order notation. One may take a common order
parameter \(q=O((r+1)^2)\). The ordinary degree-rank chains have
length \(O(r^2)\). The multidegree constructions use at most
\(O(r)\) variables with coordinate degrees at most \(O(r)\), hence
total degree \(O(r^2)\). In particular, Lemma C.6 is used only on
these downsets, not on an arbitrary downset from its general statement.

Here is a concrete inventory for Appendix C. Write \(R(\mathbf a)\)
for the coordinate downset below \(\mathbf a\). The two-variable
splittings in Lemma 7.5, Lemma 12.1, and Section 12, Step 5, use
\(R(0,r-1)\) and \(R(r-1,r-2)\). The five-variable splitting at
the end of Lemma 7.5 uses
\[
 R(r-1,0,0,0,0),\qquad R(r-2,r-1,r-1,r-1,r-1).
\]
The three-variable errors in Section 12, Steps 1--2, use unions of
rectangles with coordinate degrees at most \(r-1\); a union of
four is handled by three binary splittings. Polarization uses the
multilinear box \(\{0,1\}^r\) and at most \(r+1\) actual variables.
All these downsets therefore have total degree \(O(r^2)\) and at
most \(\exp(O(r\log(r+1)))\) indices.

\begin{center}
\small
\begin{tabularx}{\linewidth}{@{}p{.25\linewidth}X@{}}
\toprule
Appendix C operation & Dependent construction \\
\midrule
Transitivity and pullback, C.1--C.2 &
Products, Taylor substitution, and at most one subgroup basis; all
multiindices are handled together.\\[3pt]
Cancellation, translation, dilation, roots, C.3 &
Two-copy subgroup, one quotient, and boundedly many basis operations.
The root case calls the dilation case once.\\[3pt]
Additivity, C.4 &
One three-copy subgroup, a Taylor identity, one basis, and one quotient.\\[3pt]
Multilinearization, C.5 &
One root call, at most \(2^q\) independent subgroup bases, a direct
sum, one B.4 call, and a final graph subgroup and quotient.\\[3pt]
Splitting, C.6 &
Bases for the original multiindex layers, at most \(q\) total-degree
stages, one universal group, one B.3 call, and two quotients.\\
\bottomrule
\end{tabularx}
\end{center}

The representation step left implicit in C.6 has a direct bulk
implementation. With \(g(0)=1\), write its logarithm in monomials
with coefficients in the prescribed multiindex layers. At total
degree \(a\), express all such coefficients in the fixed layer
bases and multiply by the inverse product of the corresponding
one-parameter factors. Since lower degrees have already vanished,
BCH commutators affect only degrees greater than \(a\). At most
\(q\) stages finish the representation. Its arbitrary real orbit
coefficients are not complexity inputs. The resulting universal
group has at most \(|J|d\) generators and dimension at most
\(K_q(1+|J|d)^q=K_r(1+d)^q\), as required by the bulk lemma.
There is no recursive C.5 or C.6 invocation. The subgroup bases
indexed by subsets or multiindices come from the same original
data, before taking a common product.

We now count actual dependency depth, taking the complete
Leng equidistribution theorem as one operation by
Lemma~\ref{single:equid}. Its two sequential branches have already
been paid inside that lemma. The central-target reductions in
\cite[Theorem 5.4]{lss} take at most \(q\) successive steps, and
Corollary~5.5 invokes that result once. The partition in Lemma B.3
and the converse estimate in Lemma B.5 each have one smaller-degree
child. Neither makes a new forward inverse call.

The seven stages of the sunflower construction, Lemma 8.3, are
sequential. Lemma 9.1 first uses that output and then makes a new
Lemma 7.5 and Corollary 5.5 invocation, followed by its fixed linear
algebra and the absolute-degree approximate-homomorphism primitive.
These successive uses are retained. Lemma 6.3 is then applied at most
\(r-1\) times, each time to the preceding rank's output. The
universal-group and extraction constructions of Sections 10--11
have at most \(q\) total-degree stages and a fixed acyclic helper
list. Section 12 has a fixed list of main stages using these results.

The large independent families require a different count. All subset
bases in C.5, all multiindex bases in C.6, and all factors in a
polarization expansion are constructed from the data prescribed
before their stage. If a stage has \(T\) factors and target error
\(\eta\), use error \(\eta/(10T)\) for every factor, expand the
whole expression, and make its correlation selection once. Products
of expansions use the product of their coefficient masses. They do
not pass a newly selected correlation threshold from one factor to
the next. The number of such factors is order-dependent, at most
\(\exp(q^{O(1)})\). Their direct product dimension is charged, but
this multiplies the input dimension by an order-dependent coefficient;
it does not apply a new budget power once per factor.

Multilinear expansion can be performed variable by variable with at
most \(O(q^2)\) addition levels; within a level the factors are
handled as above. The C.6 representation has its \(q\) true
sequential total-degree cancellations. Powers such as \(r!\) and
\(r^r\) affect rational coefficients and tensor output dimensions,
not the filtration degree or the number of inverse calls. The actual
multidegrees remain those in the inventory above.

A conservative bound on the resulting dependency depth is
\(C_0(q+1)^6\), where \(C_0\) is absolute. This covers the
\(O(q^3)\) sunflower/linearization blocks, their rank iteration,
and the remaining fixed nesting of degree and quotient operations.
Taking \(H_q\le\exp(q^{O(1)})\) to include the full
equidistribution operation, the complete body degree satisfies
\[
 F_r\le(2H_q)^{C_0(q+1)^6},\qquad
 \log F_r\le r^{O(1)}.
\]
This estimates the actual constructions, including their parallel
assemblies, rather than interpreting an unspecified \(O_r(1)\).
\end{proof}

\subsection{Integration of supplied derivative correlations}

We replace the last forward inverse call in Section~12 of \cite{lss}
by a recursion which retains its already constructed lower-degree
correlators. The following elementary calculation supplies the precise
recursive input. Set
\(\Delta_t^+A(n)=A(n+t)\overline{A(n)}\).

\begin{lemma}[Mixed correlations supply derivative correlations]
\label{single:cs-witness}
Let \(A\) be one-bounded on \([L]\), extended by zero, let
\(B\) be one-bounded on \([2L]\), and let \(\psi_h\),
\(h\in[L]\), be one-bounded sequences. Suppose
\[
 \frac1L\sum_{h=1}^L
 \left|\frac1L\sum_{n=1}^L A(n)B(n+h)\psi_h(n)\right|
 \ge\rho.
\]
If \(L\ge4\rho^{-2}\), there are sequences \(\Xi_t\),
\(1\le t<L\), each a product of two translated or conjugated
members of the original family, or zero, such that
\[
 \frac1L\sum_{t=1}^{L-1}
 \left|\frac1L\sum_{n\in\Z}\Delta_t^+A(n)\Xi_t(n)\right|
 \ge\rho^2/8.
\]
At least \(\rho^2L/16\) positive shifts have correlation at
least \(\rho^2/16\). If the \(\psi_h\) are degree-\(a\)
nilsequences on manifolds of dimension at most \(d\), the
\(\Xi_t\) have the same degree and dimension at most \(2d\).
No regularity in the shift parameter is required.
\end{lemma}

\begin{proof}
Choose unit phases \(c_h\) to make the sum of the inner
correlations real and at least \(\rho L^2\). With \(m=n+h\),
Cauchy--Schwarz removes \(B\) and gives
\[
 \rho^2L^4\le2L\sum_{m=1}^{2L}
 \left|\sum_{h=1}^L c_h A(m-h)\psi_h(m-h)\right|^2.
\]
Expand the square and average in its second shift \(u\).
For some \(u\in[L]\), taking absolute values termwise yields
\[
 \sum_{t\in u-[L]}
 \left|\sum_{n\in\Z}\Delta_t^+A(n)
       \psi_{u-t}(n+t)\overline{\psi_u(n)}\right|
 \ge\rho^2L^2/2.
\]
The zero-shift term is at most \(L\), so the nonzero terms
total at least \(\rho^2L^2/4\). For \(t=-v<0\), substitute
\(n=m+v\) and conjugate; the witness for \(\Delta_v^+A(m)\)
becomes \(\overline{\psi_{u+v}(m)}\psi_u(m+v)\).
At most two terms fold to one positive shift. Keep the larger
correlation, losing at most a factor two. The asserted average
and the density of shifts follow. The zero extension belongs to
\(A\), not to a nilsequence witness. Product groups and fixed
orbit translations prove the last assertion without estimating
any orbit coefficient.
\end{proof}

For the omitted small case \(L<4\rho^{-2}\), the mixed
hypothesis implies \(\E|A|^2\ge\rho^2\). Parseval on
\(\Z/L\Z\) gives a linear phase with correlation at least
\(\rho/\sqrt L\ge\rho^2/2\). Its logarithmic budget has the
required bound. This also handles small-interval alternatives below.

\begin{proposition}[Supplied derivative integration]\label{single:integration}
For \(r\ge2\), suppose a one-bounded \(f:[N]\to\mathbb C\),
extended by zero, has one-bounded degree-\((r-1)\) nilsequences
\(\theta_h\), \(h\in[N]\), satisfying
\[
 \E_{h\in[N]}
 \left|\E_{n\in[N]}\Delta_h^+f(n)\theta_h(n)\right|\ge\rho.
\]
Let \(X\ge2\) bound their dimensions, all logarithmic numerical
data, and \(\log(1/\rho)\). Then \(f\) correlates with a
degree-\(r\) nilsequence whose dimension and logarithmic numerical
and inverse-correlation budget is at most \(K_rX^{J_r}\), where
\[
 J_r\le\exp(r^{O(1)}).
\]
The witnesses may use different manifolds and arbitrary polynomial
orbits, subject only to the common budget.
\end{proposition}

\begin{proof}
Let \(F_r\) be the complete non-inverse body degree from
Lemma~\ref{lem:external-body}, enlarged by an absolute power to
include the elementary lemma and the final products. For \(r=2\),
the fixed-degree converse estimate of \cite[Lemma B.5]{lss} and
the Gowers derivative identity give a \(U^3\) lower bound with
inverse logarithm \(X^{O(1)}\). The established fixed-order
\(U^3\) inverse theorem supplies a finite absolute \(J_2\).

For \(r\ge3\), threshold the supplied correlations and begin
\cite[Lemma 6.2]{lss} immediately after its initial inverse call.
All subsequent steps accept the witnesses as data: normalization
at zero, selection of common structure constants and filtration,
selection from the fixed function partition, vertical Fourier
selection, and nilcharacter embedding. In particular, no net of
all Lipschitz functions is used. They produce precisely the
degree-rank correlation structure of Definition~6.1 of that paper.

Apply its Lemma~6.3 through the rank descent and then its
Section~12 symmetrization. Those statements require this correlation
structure, not a fresh Gowers-norm hypothesis. At the end of
Section~12, before its last use of the converse estimate, the scalar
coordinate selection gives
\[
 \E_h\left|\E_n A(n)B(n+h)\psi_h(n)\right|\ge\rho',
 \qquad \deg\psi_h\le r-2.
\]
Divide by the controlled supremum bounds. The interval and shift
ranges can be padded by zeros as in the original proof. Apply
Lemma~\ref{single:cs-witness} to obtain supplied derivative
correlations for \(A\), and apply the present assertion at
degree \(r-1\). Its output is a degree-\((r-1)\) correlator
for \(A\), exactly the object required in the last paragraph of
the original proof. Multiplication by its already known degree-\(r\)
factor, substitution, and residue-class recombination finish the
argument. The earlier converse invocation in Section~12, Step~3,
is retained and is part of \(F_r\).

If an intermediate size condition fails, \(\log N\) is bounded
by the current body budget. The supplied correlation gives a lower
bound on \(\E|f|^2\); Parseval then supplies a linear-phase
correlator with logarithmic cost \(O(\log N+\log(1/\rho))\).
Thus this alternative is within the same budget. The only recursive
integration call in the main alternative is the one at \(r-1\).
Charging the product of the before- and after-recursion body degrees
inside \(F_r\), a sufficient recurrence is
\[
 J_r\le F_r(1+J_{r-1}).
\]
Since \(\log F_r\le r^{O(1)}\), summation of logarithms
proves the asserted bound.
\end{proof}

\begin{proposition}[Single-exponential inverse budget]\label{prop:inverse}
The exponents in \eqref{eq:inverse-budget} may be chosen with
\[
 I_r\le\exp(r^{O(1)}).
\]
\end{proposition}

\begin{proof}
The \(U^2\) Fourier argument gives \(I_1=1\), and the fixed
order \(r=2\) has a finite absolute exponent. For \(r\ge3\),
the derivative identity and the degree-\((r-1)\) inverse theorem
give the input to Proposition~\ref{single:integration}, with budget
\(K_r(2+\log(1/\delta))^{I_{r-1}}\). The derivative norm and
shift-density thresholds are powers \(\delta^{C_r}\), whose
logarithms introduce coefficients only. Therefore
\[
 I_r\le(1+I_{r-1})(1+J_r).
\]
Taking logarithms and summing the bounds on \(\log(1+J_r)\)
gives \(\log(1+I_r)\le r^{O(1)}\). Both stages of the original
argument have been accounted for; the second now uses a proved
integration interface with the supplied lower-degree correlators.
\end{proof}

\section{Relative detection by cube degree}\label{sec:relative}

Throughout Sections~\ref{sec:relative}--\ref{sec:library}, let
\(\mathcal H_D\le\exp(D^{O(1)})\) be a common degree
allowance for the elementary coordinate, mask, and completion
operations. Section~\ref{sec:library} enlarges it to include the
source's nonrecursive operations. It contains no forward inverse
or relative-detector invocation.

\subsection{Effective cube completion}\label{single:completion-section}

The following statement concerns ordinary integer cubes.  It is uniform
in the coefficients of the polynomial orbit.

\begin{lemma}[Deletion of one cube vertex]\label{single:completion}
Let \(q\ge2\), let \(G_0=G_1=G\), \(G_q=1\), and let
\(Y=G/\Gamma\) have dimension \(d\) and an adapted Mal'cev basis of
complexity \(M\ge2\).  Let \(F:Y\to\mathbb C\) be one-bounded and
\(L\)-Lipschitz, where \(L\ge1\), and let \(0<\eta<1/2\).
For any vertex \(i\in\{0,1\}^q\), there are coefficients \(c_a\)
and one-bounded Lipschitz functions \(F_{a,\omega}:Y\to\mathbb C\),
for \(\omega\ne i\), such that every adapted polynomial map
\(g:\Z^n\to G\) satisfies
\begin{equation}\label{single:completion-expansion}
 \left|F(g(x+i\cdot h)\Gamma)
 -\sum_a c_a\prod_{\omega\ne i}
       F_{a,\omega}(g(x+\omega\cdot h)\Gamma)\right|\le\eta
\end{equation}
for all \(x,h_1,\ldots,h_q\in\Z^n\).
The functions and coefficients are independent of \(g\).  With
\(X=2+d+\log M+\log L+\log(1/\eta)\), they may be chosen so that
\begin{equation}\label{single:completion-budget}
 \log\!\left(2+\sum_a|c_a|\right)
 +\max_{a,\omega}\log\bigl(2+\Lip(F_{a,\omega})\bigr)
 \le K_qX^{\exp(q^C)}.
\end{equation}
Each individual factor uses the original nilmanifold of dimension
\(d\) and has degree at most \(q-1\).
\end{lemma}

\begin{proof}
Reflections and permutations of cube coordinates reduce the assertion
to deletion of the all-ones vertex.  Identify vertices with subsets of
\([q]\).  For \(S\subseteq[q]\), write \([a]_S\) for the element of
\(G^{\{0,1\}^q}\) equal to \(a\in G_{|S|}\) at vertices containing
\(S\), and equal to the identity elsewhere.  The upper-face cube group
\(\mathcal C\) consists of the ordered products
\[
 c=\prod_{S\subsetneq[q]}[a_S]_S,\qquad a_S\in G_{|S|},
\]
where the ordering extends inclusion.  The full face is absent because
\(G_q=1\).  Face commutators satisfy
\([[a]_S,[b]_T]=[[a,b]]_{S\cup T}\); commutator closure of the
filtration therefore gives this product description.  Successive
proper vertex values recover all \(a_S\), so deletion defines an
injective homomorphism
\[
 \Pi:\mathcal C\longrightarrow P:=G^{2^q-1}
\]
onto a closed rational subgroup \(H\).  Filtered polynomial maps send
integer cubes into \(\mathcal C\): their successive face coefficients
lie in \(G_{|S|}\), by the iterated-difference definition of a
polynomial map, and the last coefficient lies in \(G_q=1\).

\textit{The full lattice.}
If all proper vertex values of \(c\) belong to \(\Gamma\), triangular
recovery gives \(a_S\in\Gamma\cap G_{|S|}\); their product at the
missing vertex also belongs to \(\Gamma\).  Consequently
\begin{equation}\label{single:completion-lattice}
 \Pi(\mathcal C\cap\Gamma^{\{0,1\}^q})
       =H\cap\Gamma^{2^q-1}.
\end{equation}
Thus deletion is injective on the nilmanifold quotients as well.
There is an explicit adapted basis: if \(X_j\) has filtration weight
\(w_j\), take \(\one_{[S]}\otimes X_j\),
\(|S|\le w_j\), \(S\subsetneq[q]\), ordered first by the original
Mal'cev index \(j\).  Its dimension is at most \(2^qd\).
Copies of the same \(X_j\) commute, and at a vertex \(T\) the
\(j\)-th second-kind coordinate is exactly
\[
 z_j(T)=\sum_{\substack{S\subseteq T\\ |S|\le w_j}}t_{j,S}.
\]
Integer Boolean M\"obius inversion on proper vertices shows that its
integer-coordinate lattice is precisely
\(\mathcal C\cap\Gamma^{\{0,1\}^q}\).  Brackets go to later
original Mal'cev indices; their coefficients are inherited from the
original structure constants.

\textit{Effective inverse and external lattice gap.}
Put \(A_q=\exp(q^{C_0})\), increasing the absolute \(C_0\) below,
and let
\[
 R=\exp\!\left[K_q(1+d)^{A_q}(1+\log M)\right].
\]
We claim that completion on the image in \(Y^{2^q-1}\) is
\(R\)-Lipschitz, after enlarging \(K_q,C_0\).
In first-kind coordinates, \(\log H\) is defined, separately in
each original coordinate \(j\), by the integral equations saying
that its proper-vertex values form a Boolean polynomial of degree
at most \(w_j\).  Each equation has coefficients \(0,\pm1\) and
at most \(2^q\) terms.  M\"obius inversion also gives an integral
linear inverse for \(\Pi\) on this subspace.

All BCH and coordinate-change expressions have bracket length at most
\(q-1\).  Listing the bracket words, clearing the at most \(d^3\)
input structure-constant denominators, and bounding these polynomials
and their derivatives on bounded coordinate boxes gives bounds of
the form \(R\).  Their number and coordinate degrees are at most
\(K_q(1+d)^{A_q}\); this uses no repeated substitution indexed by
the dimension.  In particular, for every
\(\gamma\in\Gamma^{2^q-1}\), its first-kind coordinates have a
common denominator at most \(R\), independently of the magnitudes
of its integer second-kind coordinates.  If \(\gamma\notin H\),
one defining integral equation for \(\log H\) therefore has
absolute value at least \(R^{-1}\) at \(\log\gamma\).

Represent two points of the cube nilmanifold by \(a,b\) in its
coordinate fundamental box.  Their images, and
\(\Pi(b)^{-1}\Pi(a)\), have bounded coordinates of the preceding
form.  If the quotient images are \(\epsilon\)-close, some ambient
lattice element \(\gamma\) makes \(\Pi(a)\) and
\(\Pi(b)\gamma\) at most \(2\epsilon\) apart.  Coordinate and
bounded-translation estimates imply
\[
 \bigl\|\log\gamma-
       \log(\Pi(b)^{-1}\Pi(a))\bigr\|\le R\epsilon.
\]
For \(\epsilon\) below an explicitly bounded reciprocal power of
\(R\), the external lattice gap forces \(\gamma\in H\).
Equation~\eqref{single:completion-lattice} then makes this an actual
cube-lattice correction.  The linear inverse and coordinate estimates
bound the intrinsic quotient distance by \(R^{O(1)}\epsilon\).
For larger \(\epsilon\), use the bounded diameter of the cube
fundamental box and divide by that explicit threshold.  Evaluation
at the missing vertex is a homomorphism with a \(0,1\) Lie matrix,
so it has the same type of metric bound.  Absorbing these fixed powers
proves the claimed global completion bound.  This argument supplies
the effective inverse; qualitative compactness alone would not do so.

\textit{Uniform tensor interpolation.}
Let \(m=2^q-1\), and let \(K\subseteq Y^m\) be the compact image
of deletion.  The function \(F\circ\mathrm{completion}\) on \(K\)
is one-bounded and \(LR\)-Lipschitz in the maximum product metric.
Extend its real and imaginary parts by the McShane formula and project
radially onto the closed unit disk.  The resulting function
\(\widetilde F:Y^m\to\mathbb C\) is one-bounded and
\(C LR\)-Lipschitz.

Choose a \(\delta/2\)-net \(z_1,\ldots,z_J\) in \(Y\), where
\(\delta=\eta/(10CLR)\).  Mal'cev coordinate boxes give
\[
 \log J\le K_q(1+d)^{A_q}
             (1+\log M+\log(1/\delta)).
\]
Set \(b_j(y)=\max(0,1-d_Y(y,z_j)/\delta)\) and
\(\phi_j=b_j/\sum_l b_l\).  The denominator is at least \(1/2\),
so \(0\le\phi_j\le1\), \(\sum_j\phi_j=1\), and
\(\Lip(\phi_j)\le CJ/\delta\).
For \(a=(a_1,\ldots,a_m)\), take
\(c_a=\widetilde F(z_{a_1},\ldots,z_{a_m})\) and use
\(\phi_{a_j}\) at the \(j\)-th retained vertex.  The tensor
partition is a convex combination of unit-disk coefficients, so the
\emph{entire sum} is one-bounded on all of \(Y^m\).  A nonzero
term has its center tuple within
\(\delta\) of the input tuple.  Its uniform approximation error is
therefore at most \(\eta\).  Also \(|c_a|\le1\) and
\(\sum_a|c_a|\le J^m\).  The displayed net estimate and
\(\log R\) prove~\eqref{single:completion-budget}, after increasing
the absolute exponent.  Restriction to \(K\) proves
\eqref{single:completion-expansion}.
\end{proof}

For cyclic variables evaluated in integer representatives, the lemma is
used only after fixing a sufficiently short box tuple and its wrap
pattern.  The resulting virtual integer cube has vertex representatives
\(u_\omega=y_\omega-a_\omega P\), with fixed integer wrap vectors
\(a_\omega\).  Apply the lemma to \(g'(y)=g(y-a_iP)\).
A retained factor is then evaluated at
\(g(u_\omega+(a_\omega-a_i)P)\), a fixed translation of its
polynomial orbit.  Such a translation changes neither the nilmanifold
nor the Lipschitz function.  Box selection and boundary errors are
accounted for separately; no dependence on \(P\) or on orbit
coefficients enters Lemma~\ref{single:completion}.

\subsection{One inverse call at each cube order}

\begin{proposition}[Relative detection by degree descent]\label{single:relative}
The conclusion of \(\M\), Lemma 5.2, admits a budget exponent
\(R_q\) satisfying, after an absolute enlargement of the
nonrecursive allowance \(\mathcal H_D\),
\begin{equation}\label{single:relative-rec}
 R_1\le \mathcal H_D,\qquad
 R_q\le \mathcal H_D(I_{q-1}+1)(R_{q-1}+1)\quad(2\le q\le D).
\end{equation}
The output remains warm. Required errors in the bounded-grid and
spatial marginal estimates have logarithmic budgets of the same
type and are charged to the cold record. In particular,
\(R_q\le\exp(D^{O(1)})\).
\end{proposition}

\begin{proof}
For \(q=1\), independent group sites give the detector mean, so
the niltest is constant. Assume henceforth that \(q\ge2\).
Use a product of cyclic groups
\(H=\prod_a\Z/P_a\Z\), where \(P_a\) is prime and
\[
 4qH_a\le P_a\le8qH_a.
\]
Thus \(P_a>q\), and box/group normalizations cost at most
\(K_q^n\). The proof of \(\M\), Lemma 5.1, uses representative
wrap cells and leading symbols, so this change of comparable
moduli preserves its bounded-grid and marginal conclusions.

Put the detector at vertex \(0\). Retain the inverse-free
densification on pp.~37--38 of \(\M\): it makes every other
factor one-bounded, leaving \(|F_0|\le B_{\rm cap}\nu\) and a correlation
\(\delta_0\) with
\[
 \log(1/\delta_0)\le K_q(1+\log B_{\rm cap}+\log(1/\delta)).
\]
The repeated Cauchy--Schwarz and interpolation steps here act
affinely on these logarithms. Their number changes \(K_q\).

Choose the opposite vertex \(i=\mathbf1\). Its conditional
dual \(D_i\) obeys \(|D_i|\le2B_{\rm cap}\), by pairwise uniformity
of distinct group-cube sites. The current correlation gives
\(\E|D_i|^2\ge\delta_0^2\). Insert \(D_i\) at that vertex
and use Gowers--Cauchy--Schwarz and \(\|F_0\|_{U^q}\le2B_{\rm cap}\),
the latter following from the majorant moments. Consequently
\[
 \|D_i/(2B_{\rm cap})\|_{U^q}\ge\delta_0^2/(4B_{\rm cap}^2).
\]
One ordinary inverse call replaces vertex \(i\) by a
one-bounded degree-\((q-1)\) nilsequence \(T\), retaining
correlation \(\rho\). Its dimension, logarithmic complexity,
and \(\log(1/\rho)\) have degree
\(\mathcal H_D(I_{q-1}+1)\) in the original growing budget, after
enlarging \(\mathcal H_D\).

\textit{Localization and the missing mask.}
Partition every representative coordinate at every vertex into
discrete intervals of representative diameter less than
\(P_a/(100q)\), allowing singletons, with boundary cutoffs
placed halfway between adjacent grid points. There are at most \(K_q^n\) tuples of
cells. One retains correlation \(\rho_1\ge\rho/K_q^n\).
Keep their indicators as one-site factors. For proper vertices
define the virtual integer cube
\[
 y_\omega=u_0+\sum_a\omega_a(u_{e_a}-u_0).
\]
The short cells fix
\(y_\omega-u_\omega=a_\omega P\) for every proper
\(\omega\); their proper-subcube faces determine these offsets
without the missing vertex. The selected full tuple also fixes
\(y_i-u_i=a_iP\), with all offsets bounded in terms of \(q\).

Smooth the \(i\)-th indicator from inside and multiply it into
\(T\). This is essential: an arbitrary indicator left at \(i\)
would prevent deleting that vertex. Its error is bounded by
\[
 \E|F_0(u_0)|\one_{\partial I_i}(u_i)
 \le2B_{\rm cap}\,\mathbb P(u_i\in\partial I_i),
\]
again by pairwise uniformity. A relative transition width
\(\theta\) with \(B_{\rm cap} n\theta\ll\rho_1\) suffices. For
sub-grid widths the half-grid endpoints permit exact agreement
on the discrete grid. Realize the selector on circles of period
\(K_qP_a\), large enough that its only supported lift in the
possible range of \(y_i-a_iP\) is the selected real cell.
It adds \(n\) dimensions and has logarithmic Lipschitz cost
\(O_q(\log(1/\theta))\); no \(\log P_a\) cost is needed.

Apply Lemma~\ref{single:completion} to this modified nilsequence
on the common orbit \(t\mapsto g(t-a_iP)\). On the other
vertices its arguments are the fixed translates
\(u_\omega+(a_\omega-a_i)P\), which do not change orbit
complexity. The identity remains valid on the entire region
defined by the other vertex masks: if the selector contributes,
\(y_i-a_iP\) is in its chosen real cell and is congruent to
\(u_i\), so it equals \(u_i\); conversely the selected cells
fix the required offset whenever \(u_i\) lies there. Otherwise
both masked values vanish. Thus deleting \(i\) introduces no
extra scoring region, and uses no periodicity of \(g\).

Request uniform completion error \(\eta\ll\rho_1/B_{\rm cap}\).
Its cube-average error is at most \(2B_{\rm cap}\eta\). If \(M\)
is the coefficient mass of the completion expansion, one term
retains correlation \(\sigma\ge\rho_1/(2M)\). All its
proper-vertex factors are one-bounded; only the detector is
sparse. Every new one-site nilsequence has the dimension of the
single inverse-output group, including the selector circles.
The possibly large number of summands contributes \(\log M\),
not a product of that many nilmanifolds.

\textit{A lower-order cube norm.}
For \(t=1,\ldots,q-1\), group the remaining nonzero vertices
by Hamming weight \(t\). The direction
\(v_t=(t,-1,\ldots,-1)\) annihilates their linear forms,
whereas its value at the detector form is \(t\), invertible
on \(H\). Cauchy--Schwarz along these \(q-1\) directions
therefore gives, for arbitrary \(F\) and one-bounded \(f_\omega\),
\[
 \left|\E_{x,h}F(x)
   \prod_{\omega\ne0,\mathbf1}f_\omega(x+\omega\cdot h)\right|
 \le\|F\|_{U^{q-1}(H)}.
\]
Indeed the final detector arguments are
\(x+\sum_t\epsilon_t t h'_t\); changing the independent
variables \(t h'_t\) gives the usual cube norm. No independence
of the directions \(v_t\) is required. Other factors are removed
by their pointwise bounds, so no new linear-forms condition on
\(\nu\) is used.

Write the resulting detector as \(F_0T_0\one_{I_0}\).
Its \(U^{q-1}\)-norm is at least \(\sigma\). Apply the
smaller-order relative statement to its diagonal cube, of
correlation at least \(\sigma^{2^{q-1}}\). In the original
notation retain \(J_0\) and replace only the arbitrary bounded
input \(P_0\) by \(P_0T_0\one_{I_0}\). The returned
degree-\((q-2)\) niltest multiplied by \(T_0\) is an
ordinary degree-\((q-1)\) niltest.

Finally smooth the root mask at the already determined correlation
scale. Lemma 5.1's spatial/chart marginal bound gives the error
\[
 \E|F_0|\,|\one_{I_0}-W_{I_0}|
 \le B_{\rm cap}\E\nu\one_{\partial I_0}
 \le B_{\rm cap}\bigl(O_q(n\theta)+\epsilon_{\rm cold}\bigr).
\]
Half-grid endpoints again cover sub-grid widths. This uses no
pointwise bound for \(\nu\). Put this degree-one mask into
the final niltest. The case \(q=1\) is the direct mean argument.

To return to the atom in \(\M\), Theorem 3.4, use its unchanged
chart-measure normalization on pp.~39--40. In that notation,
\(\chi=1\) on the support of \(J_0\) and \(a_\chi/m_j\) is warm,
so the returned atom is \(J_0\) times the degree-\((q-1)\) niltest.
All output choices above are warm; only the marginal and grid
comparison errors enter the late cold requests.

There is one ordinary inverse call, one completion operation,
and one smaller-order relative call. Products and all required
errors have the same ordered budgets, proving
\eqref{single:relative-rec}. Taking logarithms of that recurrence
and using Proposition~\ref{prop:inverse} gives
\(\log(R_q+1)\le D^{O(1)}\), as claimed.
\end{proof}


\section{The source-local numerical schedule}\label{sec:library}

\begin{proposition}[Local polynomial degrees]\label{prop:local}
For the construction above and the imported analytic framework, the
fresh absolute rank is at most \(K_kp^F\), independently of the
incoming path dimension. A common degree \(B\) bounds all nested
structural path-dimension polynomials, the width-increment
polynomials after eliminating the intermediate numerical records,
and the root logarithmic cost and side-requirement polynomials.
One may choose
\begin{equation}\label{eq:FB}
 1\le F\le B,\qquad 2\le B\le\exp(D^{O(1)}).
\end{equation}
At weight \(i\), the width increment may use only incoming
\(Q_j\) with \(j>i\). Root requirements may use all widths.
\end{proposition}

\begin{proof}
There are two separate estimates: the degree of a nonrecursive
operation in its growing inputs, and the number of successive
substitutions of such estimates. We record both.

\textit{Nonrecursive operations.}
Include every growing dimension and coordinate count in a common
budget \(x\). Joint rational operations and bounded-degree
coefficient arrays satisfy Lemma~\ref{lem:bulk}, with an order
parameter polynomial in \(D\). A degree-\(D\) polynomial in
\(n\) growing variables has at most \((n+1)^D\) coefficients.
In contrast, the Boolean jet of degree \(h\) on a cube of order
\(q\) has at most
\[
 j(h,q)=\sum_{a=0}^{\min(h,q)}\binom qa\le2^q
\]
coordinates per axis. Its fixed order-dependent size is not a new
inverse-theorem order or a new determining weight.

The remaining elementary analytic bounds also have uniform
polynomial degrees. The lattice-extraction bound in \(\M\), Lemma B.2,
has the explicit form
\[
 C_* =\delta^{-1}\bigl(C(n+m)R\bigr)^{C(n+m)};
\]
its logarithm is polynomial in the dimensions and the logarithmic
requests. The chart matrices in Appendix F have logarithmic norm
\(O(d^3)\). For the real sublevel estimate in Appendix G, the
reciprocal exponent is \(Nd\); the one-axis minor uses at most
\(j(h,q)h(q+1)\) variables and degree at most \(j(h,q)(h-1)\).
Its order-dependence is therefore at most \(\exp(D^{O(1)})\).
The simultaneous Jacobian is block diagonal. Applying the estimates
axiswise and summing errors adds the logarithm of the number of
axes; it does not multiply the polynomial degrees of all minors.
The many-multiples induction in Lemma G.3 has a recurrence
\(C_a\le K(4C_{a-1}+10)\), with an absolute \(K\), and hence
exponential dependence on the distinguished monomial degree.

The modular estimates in Appendix H use differencing loss
\(2^{1-h}\) and a block count of size
\(K_h(C+m+1)\). Those block counts are structural and are fixed
before numerical accuracy. Its exceptional-prime argument tests
individual prime powers and product witnesses. The large integer
depending only on \(h\) enters as an additive constant in a
logarithmic height, not as a power of \(x\). Likewise flag
truncation in Section 7.4 has
\(\log Q_d\le 2\log T\log Z/\log2\), a product of two
growing logarithms. No step takes the product of all moduli below
an accuracy cutoff. If a growing Fourier derivative order is
needed, choose the initially free smooth cutoffs as standard
Gevrey bumps with derivative constants bounded by
\(C^{a+1}(a!)^2\). Their logarithms are polynomial in \(a\);
scaling adds \(a\) times the scale logarithm.

The selected-plane construction has an explicit structural
dimension bound as well. For \(\ell\) parameter directions and
output degrees \(d_i\le D\), its equation and derivative counts are
\[
 J=\sum_i\left[\binom{\ell+d_i}{d_i}-1\right],\qquad
 J'=\sum_i\binom{\ell+d_i-1}{d_i-1}.
\]
The moment and derivative failure estimates in the proof of
Lemma 6.1, pp.~48--49, are satisfied by a free dimension
\(K_D(\ell+J+J'+1+\log(1/c_*))\).
Thus it is polynomial in the structural inputs with degree
\(O(D)\). Hensel lifting has an exact normalized count at every
prime-power depth and does not update a global budget at each depth.

For the degree-one comparison, the dependence on the fixed slack
must also be separated from the polynomial degree. In \(\M\),
pp.~101--107, the moments are \(O_\epsilon(x)\) and
\(\log M=O_\epsilon(x)\). The largest regularity amplification
\(M^{2q}\) has logarithm \(O_\epsilon(x^2)\), so one can choose
the regularity error \(\exp[-K_\epsilon x^2]\). The dependent
random choice then has inverse-density logarithm
\(O_\epsilon(x^2)\). Formula (C.4) gives per-round rank increase
\(O_\epsilon(x^8)\) and logarithmic width loss
\(O_\epsilon(x^2+\log(2+d))\). After \(O_\epsilon(x)\) rounds,
rank, width logarithm, and mass logarithm have powers at most
\(9,3,12\), respectively. Smoothing the Bohr tests fits within
\(K_\epsilon(2+x)^{64}\). The long concentration-peeling loop is
paid by the telescoping potential (C.8), with the allowed successor
shapes fixed beforehand. Thus changing the slack at later degrees
changes coefficients only. Positive counting in Appendix D has
at most \(s\) axis compressions and, before those compositions,
an initial logarithmic cap budget of degree \(s+2\).

These observations, together with the joint-array operations of
Sections~\ref{sec:macro}--\ref{sec:active}, give a nonrecursive
degree bound
\begin{equation}\label{eq:H}
 \mathcal H_D\le\exp(D^{O(1)}).
\end{equation}
Fourier modes, finite menus, and empirical model terms contribute
their full logarithmic counts. Selected pivot partners contribute
their structural product-group dimension. This dimension is an
input to the next operation; it is not dropped from the bound.

\textit{Parallel major refinements with sharp output records.}
The remaining apparent exponential chains are the tensor restoration
in Lemma 4.7 and the degree tuples in backward rank transfer.
Their full call trees have exponentially many nodes, but only
degree-decreasing edges carry a new cold inverse output.

To see this, use the actual arithmetic conclusion of \(\M\),
Lemma 4.6, rather than its looser polynomial-logarithmic summary:
if \(d'\) clears the substituted rational polynomial, its
denominator divides \(d'D_0\). This holds for all tensor entries
at once with a common denominator. Primitive extension and
coefficient-one extraction preserve that denominator. A local
denominator may initially vary with the productive path, so
one should not assert divisibility by a single initial integer
without first selecting it. Instead, for the logarithmic bound
\(r_t\) of the current common denominator, a safe update is
\begin{equation}\label{single:sharp-den}
 r_{t+1}\le A_D(n)\bigl(1+r_t+\log D_0+\Lambda_{>j}\bigr).
\end{equation}
Here \(A_D\) is a structural polynomial of degree at most
\(\exp(D^{O(1)})\), and \(\Lambda_{>j}\) contains permitted
higher-layer lattice data. In the native primitive coordinates
the coefficient of \(r_t\) can be taken to depend only on \(D\);
the displayed bound safely allows joint coordinate changes.
Taking a product of two denominators adds their logarithms.
There is no new power of \(r_t\).

The real conclusion of Lemma 4.7 is equally sharp: the coefficient
norm \(B_{\rm slow}\) is multiplied by a dimension-controlled factor,
independently of the I.1 output complexity. The factors in (4.6)
are canceled by the final normalized \(z^\alpha\) multiplication.
Thus its logarithmic norm has an affine update of the same kind
as \eqref{single:sharp-den}. All scalar entries are processed
jointly. The number \(L\) of degree stages or degree tuples
depends only on \(D\), not on the structural dimension. Even
allowing \(L\le\exp(D^{O(1)})\), iterating these affine
updates gives a structural power at most
\[
 L\deg A_D\le\exp(D^{O(1)}).
\]
It does not give \((\deg A_D)^L\) as a power of the
precision variables. Fixed-order multipliers may be absorbed
into \(K_D\).

In Lemma 4.7, pp.~32--33, the parent first calls I.1 and fixes
\(K'\), its threshold, and its preliminary height data. Every
extra-sheet-degree child is then invoked \emph{starting with this
same \(K'\)}. A child may internally restrict to \(K''\),
but extends its resulting identity back to its input \(K'\).
The next sibling receives the sharp slow/rational identity and
its common denominator, not that child's \(K''\), cold height,
or inverse-output complexity. The forward bound and denominator
bound hold on all the original productive path substitutions,
so subtraction does not require passing a newly selected internal
path family to the next sibling. Forecast the sharp affine
envelope for every sibling, use that common input bound, and take
the maximum of their late requirements.

A parent-to-child edge really does carry the parent's newly
frozen sides and \(K'\) description, and is counted. Its
weighted degree decreases by at least one. Put
\(I=\max_{r\le D+2}I_r\) and \(R=\max_{q\le D}R_q\).
One I.1 block has a degree bound
\[
 V_D\le[2\mathcal H_D(I+R+1)]^{D^{C_0}}
\]
for an absolute \(C_0\). Here the block consists of the original
translation-group realization, the F.2 Property~2 and Appendix~H
detector, the B.1 degree descent, the I.3--I.4 algebraic
reconstruction, and the joint scalar components. Their genuine
dependency depth is polynomial in \(D\); they do not call
Lemma 4.7 or backward rank transfer. Thus this definition has no
dependence on the refinement being estimated. If \(E_d\) denotes
the resulting degree at weighted degree \(d\), a sufficient
recurrence is therefore
\begin{equation}\label{single:major-depth}
 E_0\le \mathcal H_D,\qquad
 E_d\le \mathcal H_DV_D\bigl(1+\max_{e<d}E_e\bigr),
 \qquad E_d\le(2\mathcal H_DV_D)^{O(d)}.
\end{equation}
The full tensor tree remains paid through its joint counts and
affine coefficient envelope; it is not treated as a serial list
of new I.1 precision powers.

At an ancestor in pp.~33--35, at most \((D+1)^D\) degree tuples
are processed. Exact separation identifies the new slow tensor
with the old local slow tensor, and the new rational tensor with
the old rational tensor, \emph{before} multiplying by deviation
or raw monomials. This prevents the unrestricted coefficients of
later determining maps from entering the next residual budget.
The sharp records again satisfy \eqref{single:sharp-den}; higher
widths enter through at most \(D\sum_{i>j}\log((w_i\epsilon_i)^{-1})\).
No cold I.1 output from one tuple is needed for the next.

Degree-zero pairs can require a productive-fraction selection.
For \(r\) tensor entries, denominator bound \(Q\), and bounded
difference radius \(C\), their number has logarithm at most
\[
 O\bigl(r^2\log(2+Q)+r\log(2+C)\bigr)
\]
in fixed higher-layer primitive coordinates. A translate of the
bounded region has the same bound. These losses add over the
order-bounded tuple list; they depend on its sharp affine envelope,
not on earlier I.1 output complexities. Forecast their total once
and use a common input fraction for every tuple. Internal further
selections remain inside the already pruned modular-good branch;
they do not require a new \(D_0\) based on the smaller fraction.
These selections prove existence of the split and, as stated at the end
of p.~35 of \(\M\), are not propagated as ancestor coefficient
or denominator costs. All late side requirements are nevertheless
retained, by taking their maximum and passing the full cold record
to lower layers. There are at most \(D\) genuine ancestor passages.

\textit{The remaining schedule.}
Appendix B has one smaller-degree child and at most \(D\) grade
corrections. Appendix C has two successive modeling phases and
one smaller-degree comparison; its second phase is uniform over
the first list. Appendix E has at most \(D\) fixed-map budget
levels. Appendices G--H have degree-index product-block estimates
and the explicit fourteen-stage schedule. I.6 has \(D\) degree
descents and \(D\) grade corrections, with pivot partners requested
together. I.15 has one smaller-old-weight child. Extraction and
the added geometry have at most polynomially many layer and
degree descents. Rank cuts use joint row systems and additive
costs, as in Section 4.6 and I.11--I.12 of \(\M\).

The four-stage numerical order of \(\M\), Section 10.3, is
unchanged. The full record \(U_j\) includes current preliminary
and late costs and passes them to lower layers. It does not feed
them back into that layer's width increment. Eliminating these
records takes \(D\) further descending substitutions. Structural
path-variable counts and precision-dependent model dimensions
are included in their respective inputs throughout.

After flattening the fixed helper graph, the genuine nonlinear
depth is \(D^{O(1)}\). The order-bounded enumerations just
analyzed contribute joint structural degrees and affine losses
within \(\mathcal H_D\le\exp(D^{O(1)})\). Proposition~\ref{single:relative}
replaces the old cube-vertex chain. A sufficient common bound is
\begin{equation}\label{eq:local-bound}
 F,B\le[2\mathcal H_D(I+R+1)]^{D^{O(1)}}
       \le\exp(D^{O(1)}).
\end{equation}
This includes all nested structural path-dimension polynomials.
The fresh-rank independence is the bottom absolute rule,
\(\M\), Lemma 2.2, obtained from additive relative rank with
zero old rank. Width increments retain the strict higher-layer
dependence of the original contracts. This proves the proposition.
\end{proof}


\section{Finite-horizon density iteration}\label{sec:iteration}
We give the complete final calculation, including a separate allowance
for the structural path-variable counts.

\begin{proof}[Proof of Theorem~\ref{thm:main}]
Let \(F,B\) be the exponents in Proposition~\ref{prop:local}, enlarged
so that \(1\le F\le B\) and \(B\ge2\). The bottom rule creates at
most \(K_kp^F\) fresh absolute coordinates per round, independently
of the incoming path dimension. All dimensions are forecast before
any numerical precision is chosen.

\textit{Dimensions and cuts.}
In one preparation round the original dimension ledger of \(\M\),
Proposition 10.2, is
\[
 H_D=d_D,\qquad H_i=d_i+H_{i+1}^2,\qquad
 d_i^{\rm out}\le d_i+d_0+H_{i+1}^2,
\]
with \(H_{D+1}=0\). A cut at \(i\) decreases \(\dim W_i\);
only exports from above can replenish this dimension. Thus \(H_i\)
bounds its cut count within the round, including restarted tests.
Over \(T\le K_kp\) rounds starting from no slots, sufficient
exponents for the dimension envelopes are
\[
 a_D=F+1,\qquad a_i=1+2a_{i+1},\qquad
 a_i=(F+2)2^{D-i}-1.
\]
Indeed the additional factor \(p\) pays the number of rounds;
descending induction bounds both the incoming dimensions and the
within-round \(H_i\). Therefore
\[
 a_0=2^D(F+3)
\]
bounds all coordinate counts and the cumulative number of cuts as
\(K_kp^{a_0}\). The structural path-column constructions are
polynomials in these counts of degree at most \(B\), by the local
budget proposition. We consequently use the common structural forecast
\begin{equation}\label{eq:structural-a}
 d_*\le K_kp^a,\qquad a=B\,2^D(F+3).
\end{equation}
The extra factor \(B\) explicitly pays for these path variables and
coefficient arrays formed from width-independent structural variables.
Precision-dependent model and partner dimensions are included in the
local numerical operations, not in \(d_*\). Thus this structural
forecast introduces no dependence on a width.

\textit{Widths and root requirements.}
The strict higher-layer dependence of the local width polynomials gives
\[
 Q_i^{\rm out}\le Q_i+
 K_k\left(2+p+d_*+\sum_{j>i}Q_j\right)^B.
\]
Starting with trivial widths, let \(K_kp^{v_i}\) bound the width
logarithm at \(i\) throughout the proposed run. We may take
\begin{equation}\label{eq:v-rec}
 v_D=1+Ba,\qquad v_i=1+Bv_{i+1}\quad(i<D).
\end{equation}
Here \(v_{i+1}\ge a\), and each added one again pays the sum over
rounds. Hence \(v_1\le(2B)^Da\). Evaluating the root-cost and
side-requirement polynomials and summing over the run gives total
logarithmic cost at most \(K_kp^A\), where
\begin{equation}\label{eq:A}
 A=1+Bv_1\le(2B)^{D+1}a.
\end{equation}
No current or lower width is substituted into its own increment.

Since \(F\le B\) and \(\log B\le D^{O(1)}\),
\eqref{eq:structural-a}--\eqref{eq:A} imply
\begin{align*}
 \log A
 &\le(D+2)\log(2B)+D\log2+\log(F+3)\\
 &\le D^{O(1)}=k^{O(1)}.
\end{align*}
Increasing one absolute constant \(C\), uniformly also for the finitely
many small orders, we obtain
\begin{equation}\label{eq:single-A}
 A\le\exp(k^C).
\end{equation}

\textit{The contradiction and inversion.}
Use one fixed starting budget \(p\), enlarging its order-dependent
coefficient if necessary so all temporary density thresholds exceed
\(e^{-K_kp}\). Reserve \(T\le K_kp\) rounds so that the target
\((\alpha/2)(1+\eta_k)^T\) reaches one, where \(\eta_k>0\) is
the retained multiplicative gain. The cost bound above sustains all
rounds whenever
\[
 \log N\ge K_k\bigl(2+\log(1/\alpha)\bigr)^A.
\]
The original certificate is impossible at target at least one, since
its positive payoff is pointwise at most its enlarged negative
baseline. This is the finite-horizon argument of \(\M\),
Proposition 10.4; endpoint and side-rounding costs were included above.

Taking \(\alpha=r_k(N)/N>0\) and using the contrapositive yields
\[
 \alpha< e^2\exp\!\left[-K_k^{-1/A}(\log N)^{1/A}\right].
\]
The case \(\alpha=0\) is immediate. By \eqref{eq:single-A},
\(1/A\ge\exp(-k^C)\). Since \(\log N>1\) for \(N\ge3\),
decreasing this exponent and renaming the positive order-dependent
constants proves \eqref{eq:main}.
\end{proof}

\begin{thebibliography}{9}
\small
\bibitem{openai}
OpenAI, \emph{Quasipolynomial Bounds for Arithmetic Progressions},
September 23, 2026, 198 pp., OpenAI/math, family 159.
\href{https://github.com/openai/math/blob/main/preprints/Quasipolynomial-Bounds-for-Arithmetic-Progressions-September-23-2026/paper.pdf}{Source manuscript}.
\bibitem{lss}
J. Leng, A. Sah, and M. Sawhney,
\emph{Quasipolynomial bounds on the inverse theorem for the Gowers
\(U^{s+1}[N]\)-norm},
\href{https://arxiv.org/abs/2402.17994v3}{arXiv:2402.17994v3}.
\bibitem{leng}
J. Leng, \emph{Efficient Equidistribution of Nilsequences},
\href{https://arxiv.org/abs/2312.10772v5}{arXiv:2312.10772v5}.
\end{thebibliography}
\end{document}
